\begin{claim}{B.11}{Proved}{The cardinal-controlled Kan--Thurston statement}
Let $Y$ be a connected $\U$-small CW complex and let $\mu$ be
a cardinal belonging to $\U$. There are a $\U$-small group $G$
and an acyclic map $t:BG\to Y$ such that
$\phi=\pi_1(t):G\twoheadrightarrow\pi_1Y$ is onto, its kernel
$P$ is perfect, and every homomorphism $G\to L$ with
$|L|\leq\mu$ vanishes on $P$. Base points are chosen only to
write $\phi$; acyclicity and the conclusion about the normal
subgroup do not depend on their choice.
\end{claim}
\begin{proof}
Choose a simplicial complex $K$ and a homotopy equivalence
$h:|K|\to Y$; the existence is
\citet[Theorem~2C.5]{C_Hatcher2002}. Choose an infinite
cardinal $\lambda\in\U$ larger than $\mu$ and at least the
number of simplices of $K$. Such a cardinal exists inside the
fixed Grothendieck universe. We construct, simultaneously for
all subcomplexes $L\subseteq K$, a CW space $T(L)$ and a map
$t_L:T(L)\to|L|$, so that inclusions are actual CW inclusions,
every component of $T(L)$ is aspherical, and every inclusion
$T(L')\hookrightarrow T(L)$ is injective on the fundamental
group of each component. All choices attached to a simplex
are made once and used in every subcomplex containing it.

In dimension at most one, put $T(L)=|L|$ and $t_L=\id$.
Graphs are aspherical, and inclusions of subgraphs induce
injections on fundamental groups, as is seen from reduced
edge paths after choosing maximal forests. Suppose the
construction and its inclusion properties have been made
through dimension $n-1$, with $n\geq2$. For each $n$-simplex
$\sigma$ of $K$, the space $T(\partial\sigma)$ is connected
and aspherical. Its fundamental group embeds, by B.10, in
an acyclic nonabelian simple group $C_\sigma$ of cardinality
$\lambda$. Realize this embedding by a cellular map
\[
 a_\sigma:T(\partial\sigma)\longrightarrow BC_\sigma.
 \tag{B.11.1}
\]
The realization is possible because both spaces are connected
$K(\pi,1)$ spaces; one can send their universal covers
equivariantly into the contractible universal cover of
$BC_\sigma$, extending over cells. Let $Q_\sigma$ be its
mapping cylinder, with $T(\partial\sigma)$ at the free end.
Thus $Q_\sigma$ deformation retracts onto $BC_\sigma$,
$T(\partial\sigma)\hookrightarrow Q_\sigma$ is a CW
cofibration, and this inclusion is injective on $\pi_1$.
For a subcomplex $L$ of dimension at most $n$, put
\[
 T(L)=T(L^{(n-1)})\ \cup_{\coprod_{\sigma\in L_n}
                  T(\partial\sigma)}\
                  \coprod_{\sigma\in L_n}Q_\sigma .
 \tag{B.11.2}
\]
On the cylinder of (B.11.1), define $t_L$ by moving
$t_{\partial\sigma}(x)$ linearly inside $\sigma$ to its
barycentre as the cylinder coordinate runs from its free end
to $BC_\sigma$. Send all of $BC_\sigma$ to that barycentre.
This defines an actual continuous map extending the one on
the lower skeleton, and it commutes with all subcomplex
inclusions. All these spaces have at most $\lambda$ cells:
the bar models of the groups do, and finite products and a
countable union preserve this infinite cardinal bound.

We justify both induction assertions, including the
subcomplex injectivity which is needed later. We use the
following elementary graph-of-spaces argument. If connected
aspherical vertex spaces are joined by mapping cylinders
of aspherical edge spaces, and both edge homomorphisms
are injective, the normal form for amalgamated products and
HNN extensions embeds all vertex groups. In the universal
cover the incidence graph is the corresponding tree of
cosets. Vertex and edge inverse images are contractible.
Every finite subtree of this lifted graph of spaces is
contractible by successive attachment of a contractible
vertex space along a contractible edge cylinder. A map
from a sphere has compact image in a finite subcomplex,
and hence lies over a finite subtree. It is nullhomotopic
there. The universal cover is therefore weakly contractible,
and, being a CW space, contractible. This proves
asphericity as well as injectivity; disconnected initial
spaces are handled component by component. It is the
standard argument underlying the amalgamation statement
in \citet[\S1, before Proposition~1.1]{BaumslagDyerHeller1980}.

Apply this first to (B.11.2). Each edge group injects into
the lower-skeleton group by the induction hypothesis and
into $C_\sigma$ by its construction. It follows that
$T(L)$ is aspherical on every component and that its lower
skeleton and all $BC_\sigma$ inject on $\pi_1$.
For the stronger assertion, let $L\subseteq K$ with both
complexes of dimension at most $n$. Form first
\[
 V=T(K^{(n-1)})\cup_{T(L^{(n-1)})}T(L).
 \tag{B.11.3}
\]
This is a graph of spaces with one edge for each component
of $T(L^{(n-1)})$. Its edge injections into
$T(K^{(n-1)})$ are the dimension-$n-1$ induction assertion;
its edge injections into $T(L)$ were proved in the preceding
sentence. The argument just given embeds both vertex groups
in each resulting component of $V$. Next attach the
$Q_\sigma$ for $\sigma\in K_n\setminus L_n$. Each
attaching subgroup injects into $T(K^{(n-1)})$, which already
injects into $V$ and continues to inject at every attachment.
The final space is exactly $T(K)$: both constructions use
the same pushouts and the same chosen $Q_\sigma$. Thus
$T(L)\to T(K)$ is injective on fundamental groups.
This proves the missing subcomplex step without supposing
that an arbitrary map between aspherical spaces is injective.

Take the union over dimensions. Compactness of spheres and
homotopies in CW complexes proves the same assertions for
arbitrary subcomplexes of an infinite-dimensional $K$.
The construction does not introduce new components, so
$T(K)$ is connected when $K$ is. We now check the exact
homological property, not merely constant homology. For
every local system $M$ of abelian groups on $|L|$, induction
on dimension gives
\[
 H_*(T(L);t_L^*M)\xrightarrow{\ \cong\ }H_*(|L|;M).
 \tag{B.11.4}
\]
Indeed, in the relative pair for an $n$-simplex the target
system is constant because $\sigma$ is contractible. The
pair on the source is $(Q_\sigma,T(\partial\sigma))$.
Its long exact homology sequence, acyclicity of
$Q_\sigma\simeq BC_\sigma$, and the induction comparison
for $T(\partial\sigma)\to\partial\sigma$ identify its
relative homology with that of $(\sigma,\partial\sigma)$,
with every constant coefficient group. Excision gives the
direct sum of these relative comparisons for (B.11.2).
The long exact sequence of the skeletal pair proves
(B.11.4). Cellular chains with local coefficients commute
with the filtered union over skeleta, which proves it in
unbounded dimension as well.

Write $G=\pi_1T(K)$ and let $N$ be the normal closure in
$G$ of the embedded groups $C_\sigma$, for all simplices
of dimension at least two. The map $t_K$ kills each of
these groups. Moreover
\[
 G/N\xrightarrow{\ \cong\ }\pi_1|K|.
 \tag{B.11.5}
\]
To verify the specified map, use van Kampen in (B.11.2).
At dimension two, killing $C_\sigma$ kills precisely
the image of $\pi_1T(\partial\sigma)=\pi_1\partial\sigma$,
the usual boundary relator of that two-simplex. At a
dimension $n\geq3$, the inductive version of (B.11.5)
for $\partial\sigma$ says that its group is normally
generated by its lower-dimensional cone groups, because
$\pi_1\partial\sigma=0$. Those groups have already been
killed. No further relation in the quotient is therefore
introduced. This proves (B.11.5) by induction and then by
filtered union. In particular $\pi_1(t_K)$ is onto and
its kernel is exactly $N$. Each $C_\sigma$ is perfect;
its conjugates generate $N$, and each such conjugate lies
in $[N,N]$. Hence $N$ is perfect.

Since $T(K)$ is a connected aspherical CW space, choose
a based homotopy equivalence $b:BG\to T(K)$ inducing the
identity on $G$, and put $t=h\,t_K\,b$. Equation (B.11.4)
and \citet[Definition/Proposition~1.2]{HausmannHusemoller1979}
show that this specified map is acyclic. Its kernel is
$P=N$, transported by the chosen base-point identification.
Every homomorphism $\rho:G\to L$ with $|L|\leq\mu$
restricts trivially to each $C_\sigma$, since that group
is simple and has cardinality $\lambda>\mu$. Thus $\rho$
kills their normal closure $P$. If $K$ has dimension at
most one, there are no cone groups, $P=1$, and the same
conclusion holds. This proves all the stated cases.
\end{proof}

\begin{claim}{B.12}{Proved}{Small ind-objects modulo natural transformations see at most one-types}
Let $\mathcal H$ be the category of $\U$-small categories
whose morphisms are functors modulo zigzags of natural
transformations. Let $I$ be a nonempty $\U$-small filtered
category, and let $K:I\to\mathcal H$ be a diagram, represented
by categories $K_i$. Suppose the presheaf
\[
 H_K(X)=\colim_{i\in I}\Hom_{\mathcal H}(X,K_i)
       =\colim_i\pi_0\Fun(X,K_i)
 \tag{B.12.1}
\]
sends every functor in $\Winfty$ to a bijection. Then it
sends every functor in $\W_1$ to a bijection. The specified
unit of the reflector induces an isomorphism
\[
 (K_i)_{i\in I}\longrightarrow(\Pi_1K_i)_{i\in I}
 \quad\hbox{in }\operatorname{Ind}_{\U}(\mathcal H).
 \tag{B.12.2}
\]
In particular the small-index version of CATHOT E.17 is
settled, with its original indexing category unchanged.
\end{claim}
\begin{proof}
We first identify exactly the homotopy-invariant functor
which (B.12.1) defines. The functor
$q:\mathcal H\to\Hot$ is localization at nerve weak
equivalences: a natural transformation is a functor on
$X\times[1]$, and the two endpoint inclusions become
equal after inverting its projection to $X$. Thus the
localization of $\Cat$ factors through $\mathcal H$
and has the same universal property there.
By hypothesis $H_K$ is the pullback of a presheaf $H'$
on $\Hot$. There is a canonical identification
\[
 H'(Y)\cong\colim_i[Y,BK_i].
 \tag{B.12.3}
\]
To justify this identification rather than assume it at
each stage, use a further universe for presheaf categories.
Restriction $q^*$ is fully faithful because $q$ is a
localization. Its left adjoint, left Kan extension along
$q^{\op}$, sends the representable
$\Hom_{\mathcal H}(-,K_i)$ to the representable
$\Hom_{\Hot}(-,BK_i)$, as the Yoneda lemma verifies on
every test presheaf. This left adjoint preserves colimits.
Its unit is invertible on a presheaf already factoring
through $q$. Applying these facts to (B.12.1) proves
(B.12.3) with its natural comparison map.

Choose a small infinite cardinal $\mu$ bounding the
cardinality of the arrows of every $K_i$. Such a bound
exists because the index category and all its categories
are small in the same universe. Let
$Y=\coprod_{c\in\pi_0Y}Y_c$ be any small CW space,
choose a base point in each component, and apply B.11 to
obtain acyclic maps $t_c:BG_c\to Y_c$. Put
$T=\coprod_cBG_c$ and $t=\coprod_ct_c$. Write
$\tau_1Y=\coprod_cB\pi_1Y_c$ and let
$v:Y\to\tau_1Y$ be its canonical first truncation in
$\Hot$. The composite $vt:T\to\tau_1Y$ is represented
by the quotient homomorphisms
$G_c\twoheadrightarrow\pi_1Y_c$.

For every $i$, restriction along that quotient gives an
\emph{isomorphism of ordinary functor categories}
\[
 \Fun\left(\coprod_cB\pi_1Y_c,K_i\right)
 \xrightarrow{\ \cong\ }
 \Fun\left(\coprod_cBG_c,K_i\right).
 \tag{B.12.4}
\]
A functor from $BG_c$ chooses an object $k$ and a
homomorphism $G_c\to\Aut_{K_i}(k)$. The target group
has cardinality at most $\mu$, so this homomorphism kills
$\ker\pi_1(t_c)$ by B.11 and factors uniquely through
$\pi_1Y_c$. A natural transformation is a family of
arrows of $K_i$ commuting with the group actions. Since
the quotient is onto, commuting with all $G_c$ actions
is equivalent to commuting with all $\pi_1Y_c$ actions.
Thus objects and morphisms factor uniquely, proving
(B.12.4) even with noninvertible endomorphisms in $K_i$
and with infinitely many components. We take components
only after proving this actual isomorphism; no invalid
commutation of $\pi_0$ with an infinite product is used.

Using (B.12.3), the resulting bijection factors as
\[
 H'(\tau_1Y)\xrightarrow{H'(v)}H'(Y)
              \xrightarrow{H'(t)}H'(T).
 \tag{B.12.5}
\]
The second map is injective. Indeed B.9 says that
$[Y,BK_i]\to[T,BK_i]$ is injective for every $i$.
Filtered colimits of compatible injections of sets are
injective: an equality in the target colimit holds at
some common stage, where injectivity applies. The
composite in (B.12.5) is bijective by (B.12.4). If
$b\circ a$ is bijective and $b$ is injective, then
$a$ is bijective: injectivity is immediate, and for
$y$ choose $x$ with $ba(x)=b(y)$ and cancel $b$.
Thus $H'(v)$ is bijective. For $Y=BX$, the checked
CATHOT C.1--C.4 identify this comparison with
$H_K(\Pi_1X)\to H_K(X)$ induced by the specified unit.
It follows that $H_K$ inverts every $\W_1$ map, since
$\Pi_1$ sends those maps to equivalences of groupoids.

We finish by checking that the conclusion is the
canonical ind-isomorphism (B.12.2). The supplied
reflector $\Pi_1:\mathcal H\to\mathcal G$, where
$\mathcal G$ is the full reflected subcategory of
groupoids, induces an adjunction on small ind-categories.
This follows directly from
\[
 \Hom_{\operatorname{Ind}(\mathcal H)}((A_i),(B_j))
   =\lim_i\colim_j\Hom_{\mathcal H}(A_i,B_j)
 \tag{B.12.6}
\]
and the reflector adjunction on its individual Hom sets.
Write $L(K)=\operatorname{Ind}(\Pi_1)(K)$ and
$\eta:K\to L(K)$ for its unit. The property just proved
is orthogonality of $K$ to every $X\to\Pi_1X$.
Taking the limit over $i$ of these bijections produces
a morphism $r:L(K)\to K$ with $r\eta=\id_K$.
The object $L(K)$ has the same orthogonality property:
each target $\Pi_1K_j$ does by reflection, and filtered
colimits preserve the resulting bijections of sets.
Consequently precomposition with $\eta$ is injective
on maps $L(K)\to L(K)$. Since
$(\eta r)\eta=\eta=\id_{L(K)}\eta$, it follows that
$\eta r=\id_{L(K)}$. This proves (B.12.2).

If $Y$ is empty, (B.12.4) is the isomorphism between
two terminal functor categories, so the argument includes
that case. No uniform bound on the sizes of arbitrary
test categories $X$ was used; only the one bound on the
small family $(K_i)$ was used. An index family unbounded
in the original universe lacks this bound and is not
covered by the statement. Thus the large-index examples
retained in the CATHOT discussion do not contradict it.
\end{proof}

The construction in B.11 follows the simplicial-cone
strategy of Kan--Thurston proofs, but the cardinal and
subcomplex arguments are supplied above. We searched
for Maunder's 1981 proof, for cardinal refinements of
Kan--Thurston, and for acyclic simple envelopes. We
verified the HNN and mitosis sources used in B.10 and
the acyclic-map criterion used in B.9--B.11; we did not
obtain Maunder's full text. Consequently no assertion
about what his proof establishes is used as a premise,
and no priority claim is attached to B.11 or B.12.
\begin{claim}{B.13}{Open}{Accessibility with only the abstract derivator axioms}
For an arbitrary derivator
$\D:\Cat_\U^{\mathrm{coop}}\to\mathrm{CAT}_\V$
satisfying Der~1--4, must $\Wlc$ be generated by a
$\U$-small set of functors as a fundamental localizer?
Must its full subcategory of $\Cat_\U^{[1]}$ satisfy
the literal Loc~(9)? We likewise leave the corresponding
accessibility assertion for $\Wc$ and the literal
accessibility assertion for $\Wone$ unresolved in this
full generality. B.4 answers the first two questions
positively for every $\U$-presentable enhanced
coefficient category, and answers all three for each
concrete coefficient theory in Part C.
\end{claim}
The obstruction to extending B.4 is the missing uniform
small family of tests or a uniform bound on witnesses.
Der~1--4 and B.3 give Kan extensions and relative descent,
but do not in our proof supply such a generating family.
For the constant and locally constant classes, the first
untreated scope is an abstract derivator with no supplied
presentable enhancement or small homotopical generating
family. For $\Wone$, accessibility is also unproved here
for a general presentable coefficient category outside the
concrete families of Part C: the accessible copower argument
for equivalences has not been extended here to its
fully faithful arrows. No explicit counterexample is supplied
to either assertion. Under
Vop\v enka's principle, the set-generation questions for
$\Wlc$ and $\Wc$ are positive by
\citet[Scholie~4.2.17]{Cisinski2006}. We neither assume
that principle nor infer an unconditional counterexample
from the lack of a proof. The exact unresolved step is
to produce a single $\lambda\in\U$ and a $\U$-small
family of arrows generating the stated classes; merely
knowing their filtered-colimit closure does not do so.

\begin{claim}{B.14}{Open}{The remaining general fibration questions}
For an arbitrary Der~1--4 derivator without an enhancement, we leave
unresolved whether every Grothendieck fibration whose transports
belong to $\Wone$ universally preserves $\Wone$ under strict
pullback, and whether every small functor factors as a member of
$\Wone$ followed by such a universally $\Wone$-preserving map.
These are respectively literal Loc~(7a) and the factorization
assertion of Loc~(7F); Loc~(7c) is proved in B.7.
For the accessible localizer $W^{\mathrm{lc}}_{L_E\Sp}$ of a general
nonzero spectrum $E$, we also leave its properness, equivalently
its Loc~(7F), unevaluated beyond the cases proved in B.7.
\end{claim}
The first two gaps are exactly the steps isolated after B.7:
the direct-image argument gives cohomological invariance, while
the identification with full faithfulness and the homotopy-fibre
criterion have only been proved here with the enhancement of C.2.
The first untreated coefficient theories have no specified such
enhancement; no counterexample derivator is supplied. For the
spectrum question the exact necessary and sufficient condition is
(B.7.7), namely generation by the aspherical projections already
in the class. The loop-homology test of C.11 decides membership
but does not prove that generation statement. This evaluation is
not supplied here already for $E=K(1)$ at a fixed prime. We do
not claim that these residual questions are open in the literature.

